\documentclass[a4paper,12pt]{ctexart}
\usepackage[top=2.5cm,bottom=2.5cm,left=2.8cm,right=2.8cm]{geometry}
\usepackage{amsmath,amssymb,amsfonts}
\usepackage{graphicx}
\usepackage{booktabs}
\usepackage{multirow}
\usepackage{array}
\usepackage{longtable}
\usepackage{float}
\usepackage{xcolor}
\usepackage{hyperref}
\usepackage{caption}
\usepackage{subcaption}
\usepackage{enumitem}
\usepackage{listings}
\usepackage{makecell}

\hypersetup{colorlinks=true,linkcolor=black,citecolor=black,urlcolor=blue}
\lstset{basicstyle=\small\ttfamily,breaklines=true,frame=single}

\title{\bfseries 城市新能源公共充电网络智能规划与调度\\——基于多维需求刻画、多目标优化与分时负荷转移的建模研究}
\author{参赛队伍：\underline{\hspace{3cm}}\quad 选题：A题}
\date{2026年}

\begin{document}
\maketitle

%================================================
\begin{abstract}
本文针对城市新能源公共充电网络面临的需求时空异质性、服务覆盖不均衡、电网峰谷冲突及中长期过载风险，构建了“需求刻画—设施配置—分时调度—滚动评估”四阶段递进模型体系。

针对问题一，基于2025年10个典型区域工作日与周末24小时充电量及区域基础属性，构建综合典型日与$5\!:\!2$加权日模型，系统刻画了区域、时段、日类型三维分布规律。结果表明：区域综合日均需求排序为$4>1>6>2>5>9>3>10>8>7$，极差达$2.57$倍；全市峰值位于$17\text{--}18$时$12511.57\,\mathrm{kWh}$，谷值位于$03\text{--}04$时$174.86\,\mathrm{kWh}$；$5$个区域为工作日主导型、$2$个为周末主导型，其中区域$9$周末需求较工作日高$99.63\%$。为克服小样本与共线性，构建人口密度、车流量、商业POI三因素岭回归模型，经留一交叉验证优选正则化参数$\lambda^{*}=10$，得到标准化回归方程，LOOCV-RMSE为$3426.49\,\mathrm{kWh/day}$。在区域属性短期稳定假设下，以2025年真实需求为基准，预测2026年基准情景（年增长率$15\%$）全市日均需求为$158708.05\,\mathrm{kWh/day}$，较2025年增加$20701.05\,\mathrm{kWh/day}$。

针对问题二，以问题一$15\%$基准情景预测车次需求为输入，构建以新增成本最小、全市面积加权覆盖率最大、新增接入容量压力平方和最小为目标的多目标整数规划模型，设置区域覆盖率$\ge 90\%$、服务能力充足、快充结构不退化等硬约束。采用NSGA-II生成Pareto前沿，经最小快充结构修复与$\pm1$邻域支配扫描消除冗余解，再经熵权-TOPSIS择优。对比压力方差模型$M\text{-}B$发现，压力平方和模型$M\text{-}A$以低$10.7\%$的成本获得更低的平均接入压力。最终推荐方案为新增$634$台快充桩与$416$台慢充桩、合计$1050$台，总投资$4136.80$万元，全市面积加权覆盖率提升至$90.61\%$且各区域均达标，平均新增接入压力率为$0.0367$。

针对问题三，构建“峰谷差最小—负荷方差最小”字典序两阶段优化模型，将高峰时段$20\%$负荷转移至低谷$7$小时，在电量守恒与逐时容量约束下采用\texttt{linprog-HiGHS}与\texttt{CVXPY/CLARABEL}求解$20$个“区域$\times$日期类型”子问题。结果显示：工作日平均峰谷差改善$30.78\%$、方差改善$62.33\%$，全市工作日峰谷差由$14057\,\mathrm{kW}$降至$9669\,\mathrm{kW}$（改善$31.2\%$）；周末平均峰谷差改善$16.01\%$、方差改善$45.56\%$，全市周末峰谷差由$11896\,\mathrm{kW}$降至$10484\,\mathrm{kW}$。调度后最大负荷率不超过$52.30\%$、最小安全裕度$614.0\,\mathrm{kW}$，原$3$个超$2100\,\mathrm{kW}$的高负荷小时全部降至阈值以下。

针对问题四，以问题二已落地网络为初始状态，按年增长率$15\%$投影2026—2028年风险。服务侧最小保障率在2028年仍达$614.0\%$（区域$10$），无服务缺口，滚动补桩MILP最优追加量为$0$；电网侧区域$9$在$10\text{--}11$时平段率先过载，调度后峰时风险率在2026—2028年分别为$108.65\%$、$124.94\%$、$143.68\%$，对应最低技术容量缺口$0.175$、$0.506$、$0.885\,\mathrm{MW}$，需实施平段需求响应与分期配网增容。模型通过敏感性检验与鲁棒性分析，具有良好的可推广性。

\textbf{关键词：}充电需求预测；岭回归；NSGA-II；多目标整数规划；熵权-TOPSIS；分时电价；负荷转移；滚动优化
\end{abstract}

\tableofcontents
\newpage

%================================================
\section{问题重述}
在“双碳”目标与新型电力系统建设背景下，新能源汽车保有量持续快速增长，城市公共充电网络面临需求时空分布不均、站桩布局不合理、电网峰谷负荷冲突及区域供电容量紧张等多重挑战。科学开展充电需求预测、站点优化选址、充放电动态调度与安全运行分析，对提升充电服务质量、保障电网安全稳定具有重要意义。

某市主城区划分为$10$个典型功能区域，涵盖老城核心区、城市新区、工业区、城郊过渡区等，各区域在人口密度、车流量、商业繁华程度、既有充电桩数量及电网容量等方面差异显著。不合理的充电行为将加剧电网峰谷差，甚至引发过载风险。附件$1$—$5$提供了2025年各区域基础属性、$24$小时分时段充电需求、电网安全容量等数据。

需建立数学模型解决以下问题：

\textbf{问题一} 基于区域基础数据与分时段充电负荷数据，从区域空间、$24$小时时段、工作日/周末三维刻画充电需求分布规律与差异特征；综合人口密度、车流量、商业POI等因素构建多因素充电需求预测模型，预测各区域未来短期日均充电需求（$\mathrm{kWh/day}$）。

\textbf{问题二} 在保持既有充电桩基础上，新增一批快充桩与慢充桩，以投资成本最小、服务覆盖率最高、电网负荷均衡性最优为目标，结合问题一预测需求、既有桩数与电网总容量，确定各区域新增快、慢充桩数量。

\textbf{问题三} 拟采用“峰、平、谷”三段式分时电价引导用户将高峰$20\%$充电负荷转移至低谷时段。基于附件$3$负荷数据计算调度前各区域峰谷差，将转移电量在低谷时段优化分配，给出调度后$24$小时负荷曲线并评估调度效果。

\textbf{问题四} 假设未来$3$年新能源汽车保有量每年增长$15\%$，评估服务覆盖率下降或电网过载风险，并提出优化调整建议。

%================================================
\section{问题分析}
本文将四个问题视为“需求刻画—设施配置—运行调度—中长期风险”递进链条，整体技术路线如图1所示。

%% 【此处粘贴：四问题递进技术路线图，含数据输入、模型衔接与输出关系】

\subsection{问题一的分析}
问题一要求揭示充电需求的分布异质性并完成短期预测。由于需求同时受区域功能、人口、交通、商业活跃度及用户出行习惯影响，需从空间、时段、日类型三维寻找规律；预测部分以区域日均需求为因变量，人口密度、车流量、商业POI为解释变量建立多因素模型。考虑到仅$10$个区域样本且解释变量间存在共线性，需采用正则化方法提升稳定性，并结合保有量增长情景修正预测结果，为后续规划提供需求输入。

\subsection{问题二的分析}
问题二要求在既有网络上确定各区域新增快、慢充桩数量。快、慢充桩在成本、服务能力、接入功率上差异显著，且各区域预测需求、既有覆盖与电网容量不同，不能按需求简单分配。首先需利用问题一预测车次需求与既有服务能力确定扩容下限；再利用既有覆盖面积反推单桩有效覆盖效率，量化新增桩对覆盖率的提升。需在满足区域覆盖率、服务能力、快充结构不退化等硬约束下，以新增成本最小、面积加权覆盖率最大、新增接入容量压力最小为目标建立多目标整数规划，通过Pareto优化择优实现经济性、服务性与电网安全性的兼顾。

\subsection{问题三的分析}
问题三要求在分时电价引导下实现削峰填谷。不同区域工作日、周末负荷曲线及逐时允许负荷差异显著，需按“区域$\times$日期类型”分别建模。首先计算调度前峰谷差；再将高峰$20\%$电量在保证全天电量守恒前提下向低谷$7$小时转移。为避免转移电量集中形成新峰值，需以调度后峰谷差最小为主目标、逐时容量为约束优化低谷分配，并在多解中进一步使低谷分布均匀；最后对比调度前后曲线、峰谷差与安全裕度评价政策效果。

\subsection{问题四的分析}
问题四需在年增长$15\%$情景下评估长期适应能力。以问题二已建网络与问题三调度结果为基线，逐年放大充电车次需求与峰时负荷，分别从服务侧（需求 vs. 服务能力）与电网侧（负荷 vs. 允许容量）识别风险。服务不足需滚动补桩，容量不足需需求响应扩展或配网增容，据此提出分区域、分年度的动态调整策略。

%================================================
\section{模型假设}
\begin{enumerate}[label=\arabic*.]
\item 一周由$5$个工作日与$2$个周末日构成，二者按$5\!:\!2$加权构造综合典型日；短期内区域人口密度、车流量、商业POI保持2025年水平——与附件短期预测缺失及城市规划稳定性相符。
\item 新增新能源汽车在各区域空间分布比例、单车公共充电频率及单次充电电量保持稳定——使保有量增长率可直接转化为充电需求增长率。
\item 现有充电桩数量是历史供给结果，可能受需求反向影响，不作为需求驱动模型的核心解释变量，避免内生性偏差。
\item 附件$3$每小时充电量按$\Delta t=1\,\mathrm{h}$转换为小时平均功率，数值不变但口径统一为$\mathrm{kW}$，用于与附件$4$逐时允许负荷比较。
\item 未提供桩位坐标、道路网络及服务半径重叠关系时，新增桩优先布设于未覆盖区域，且平均有效覆盖效率与本区域既有网络相同——以既有网络反推单桩覆盖效率。
\item 快充与慢充服务结构不低于区域现状，防止模型因慢充成本低而退化为全慢充解。
\item 分时电价峰、平、谷时段为外生政策，不作为优化变量；$23\text{--}00$时未被题设明确归类，保持原负荷不变且不参与转移。
\item 未来三年日内负荷形态及关键峰时段保持相对稳定，峰时负荷率可按年度增长系数同比例投影，用于中长期风险快速评估。
\end{enumerate}

%================================================
\section{符号说明}
\begin{longtable}{cll}
\toprule
符号 & 含义 & 单位 \\
\midrule
$i$ & 区域编号，$i=1,\ldots,10$ & --- \\
$t$ & 时段编号，$t=1,\ldots,24$ & --- \\
$d$ & 日期类型，$d\in\{\mathrm{wd},\mathrm{we}\}$ & --- \\
$E_{i,t}^{d}$ & 区域$i$在日期类型$d$、时段$t$的充电量 & $\mathrm{kWh}$ \\
$D_{i}^{\mathrm{wd}},D_{i}^{\mathrm{we}}$ & 区域$i$工作日、周末全天充电需求 & $\mathrm{kWh}$ \\
$D_{i,0},\bar{D}_{i}$ & 区域$i$综合日均充电需求（2025基准） & $\mathrm{kWh/day}$ \\
$E_{i,t}^{*}$ & 综合典型日分时充电量 & $\mathrm{kWh}$ \\
$\bar{E}_{i}^{*}$ & 区域平均小时充电量 & $\mathrm{kWh}$ \\
$R_{i}^{*}$ & 归一化极差 & --- \\
$\eta_{i}$ & 周末相对工作日差异比 & $\%$ \\
$\mathrm{PD}_{i},\mathrm{TF}_{i},\mathrm{POI}_{i}$ & 人口密度、车流量、商业POI & --- \\
$z_{ij}$ & 标准化特征 & --- \\
$\beta_{0},\boldsymbol{\beta}$ & 岭回归截距与系数 & --- \\
$\lambda$ & 岭回归正则化参数 & --- \\
$\widehat{D}_{i,\tau}$ & 未来$\tau$年预测需求 & $\mathrm{kWh/day}$ \\
$g_{\mathrm{EV}},g$ & 新能源汽车保有量增长率 & --- \\
$x_{i},y_{i}$ & 新增快充、慢充桩数量 & 台 \\
$A_{i},A_{i}^{(0)}$ & 区域总面积、既有覆盖面积 & $\mathrm{km^{2}}$ \\
$F_{i}^{(0)},S_{i}^{(0)},N_{i}^{(0)}$ & 既有快充、慢充、总桩数 & 台 \\
$Q_{i}^{\mathrm{plan}},Q_{i,k}$ & 规划期日均充电车次需求 & 车次/日 \\
$G_{i},G_{i,t}$ & 区域电网总容量、逐时允许负荷 & $\mathrm{kW}$ \\
$a_{i}$ & 单桩有效覆盖面积 & $\mathrm{km^{2}/台}$ \\
$\theta_{i}$ & 既有快充占比 & --- \\
$\nu_{i},\bar{\nu}$ & 新增接入压力率及其均值 & --- \\
$C_{i},C_{\mathrm{city}}$ & 区域覆盖率、全市面积加权覆盖率 & --- \\
$c_{f},c_{s}$ & 快、慢充单桩建设成本 & 万元/台 \\
$p_{f},p_{s}$ & 快、慢充额定功率 & $\mathrm{kW/台}$ \\
$s_{f},s_{s}$ & 快、慢充日服务能力 & 车次/(台$\cdot$日) \\
$L_{i,t}^{d},\widetilde{L}_{i,t}^{d}$ & 调度前、后平均充电负荷 & $\mathrm{kW}$ \\
$z_{i,t}^{d}$ & 低谷时段转移负荷 & $\mathrm{kW}$ \\
$M_{i}^{d}$ & 高峰转移总量 & $\mathrm{kWh/day}$ \\
$U_{i}^{d},V_{i}^{d}$ & 调度后负荷上、下界 & $\mathrm{kW}$ \\
$\Delta_{i}^{d}$ & 全天峰谷差 & $\mathrm{kW}$ \\
$\mathcal{V},\mathcal{M},\mathcal{H}$ & 谷、平、峰时段集合 & --- \\
$B_{i}^{d}$ & 低谷接纳能力 & $\mathrm{kW}$ \\
$SC_{i,k},\psi_{i,k}$ & 服务能力、服务保障率 & 车次/日, --- \\
$\rho_{i,k}^{\mathrm{tou}}$ & 调度后峰时风险率 & --- \\
$\Delta G_{i,k}^{\min}$ & 最小技术容量缺口 & $\mathrm{kW}$ \\
\bottomrule
\end{longtable}

%================================================
\section{模型的建立与求解}
\subsection{问题一：城市公共充电需求三维分布与短期预测}

\subsubsection{开篇明义}
本节基于2025年附件数据，完成三维分布刻画与多因素预测建模。技术路线为：构造综合日模型$\to$三维统计刻画$\to$三因素岭回归建模与LOOCV选型$\to$多情景短期预测，为后续规划提供需求基线。

\subsubsection{分步推导}

\textbf{1. 综合日模型构建}

设$D_{i}^{\mathrm{wd}}=\sum_{t=1}^{24}E_{i,t}^{\mathrm{wd}}$，$D_{i}^{\mathrm{we}}=\sum_{t=1}^{24}E_{i,t}^{\mathrm{we}}$，按$5\!:\!2$加权定义综合日均需求：
\begin{equation}
D_{i,0}=\bar{D}_{i}=\frac{5D_{i}^{\mathrm{wd}}+2D_{i}^{\mathrm{we}}}{7}\tag{1-1}
\end{equation}
综合典型日分时充电量为：
\begin{equation}
E_{i,t}^{*}=\frac{5E_{i,t}^{\mathrm{wd}}+2E_{i,t}^{\mathrm{we}}}{7}\tag{1-2}
\end{equation}
该模型符合一周日类型结构，消除工作日与周末样本量差异。

\textbf{2. 区域空间维度}

计算$10$区域$D_{i,0}$及排序如表1所示。

\begin{table}[H]
\centering
\caption{10区域工作日、周末与综合日均充电需求}
\begin{tabular}{ccccc}
\toprule
区域 & 工作日需求/$\mathrm{kWh}$ & 周末需求/$\mathrm{kWh}$ & 综合日均需求/$\mathrm{kWh/day}$ & 排名 \\
\midrule
$4$ & $19888$ & $15711$ & $18694.57$ & $1$ \\
$1$ & $19206$ & $14580$ & $17884.29$ & $2$ \\
$6$ & $18376$ & $13350$ & $16940.00$ & $3$ \\
$2$ & $17767$ & $13348$ & $16504.43$ & $4$ \\
$5$ & $13465$ & $19105$ & $15076.43$ & $5$ \\
$9$ & $11642$ & $23241$ & $14956.00$ & $6$ \\
$3$ & $13486$ & $11183$ & $12828.00$ & $7$ \\
$10$ & $9529$ & $9065$ & $9396.43$ & $8$ \\
$8$ & $9524$ & $5775$ & $8452.86$ & $9$ \\
$7$ & $7038$ & $7864$ & $7274.00$ & $10$ \\
\bottomrule
\end{tabular}
\end{table}

排序为$4>1>6>2>5>9>3>10>8>7$，最高与最低比值$18694.57/7274.00=2.570$，呈现显著区域分层。

%% 【此处粘贴：工作日周末全天需求对比柱状图 p1_f02】

\textbf{3. 24小时时段维度}

定义区域平均小时充电量$\bar{E}_{i}^{*}=\frac{1}{24}\sum_{t}E_{i,t}^{*}$，归一化极差：
\begin{equation}
R_{i}^{*}=\frac{\max_{t}E_{i,t}^{*}-\min_{t}E_{i,t}^{*}}{\bar{E}_{i}^{*}}\tag{1-3}
\end{equation}
峰、谷时段$t_{i,\max}=\arg\max_{t}E_{i,t}^{*}$，$t_{i,\min}=\arg\min_{t}E_{i,t}^{*}$。

\begin{table}[H]
\centering
\caption{区域日内波动与峰谷时段特征}
\begin{tabular}{cccccc}
\toprule
区域 & $R_{i}^{*}$ & 峰值/$\mathrm{kWh}$ & 峰时段 & 谷值/$\mathrm{kWh}$ & 谷时段 \\
\midrule
$1$ & $2.4675$ & $1856.43$ & $17\text{--}18$ & $17.71$ & $04\text{--}05$ \\
$2$ & $2.5126$ & $1739.43$ & $17\text{--}18$ & $11.57$ & $03\text{--}04$ \\
$3$ & $2.4405$ & $1316.29$ & $17\text{--}18$ & $11.86$ & $03\text{--}04$ \\
$4$ & $2.1819$ & $1711.29$ & $17\text{--}18$ & $11.71$ & $04\text{--}05$ \\
$5$ & $2.2095$ & $1393.71$ & $18\text{--}19$ & $5.71$ & $02\text{--}03$ \\
$6$ & $2.1175$ & $1503.14$ & $17\text{--}18$ & $8.57$ & $03\text{--}04$ \\
$7$ & $2.1319$ & $651.86$ & $12\text{--}13$ & $5.71$ & $02\text{--}03$ \\
$8$ & $2.2828$ & $812.57$ & $17\text{--}18$ & $8.57$ & $01\text{--}02$ \\
$9$ & $1.9405$ & $1215.71$ & $12\text{--}13$ & $6.43$ & $03\text{--}04$ \\
$10$ & $2.0028$ & $797.71$ & $12\text{--}13$ & $13.57$ & $04\text{--}05$ \\
\bottomrule
\end{tabular}
\end{table}

全市$E_{t}^{\mathrm{city}}=\sum_{i}E_{i,t}^{*}$峰值出现在$17\text{--}18$时$12511.57\,\mathrm{kWh}$，谷值在$03\text{--}04$时$174.86\,\mathrm{kWh}$。区域$1,2,3,4,6,8$为晚高峰型，$7,9,10$为午间高峰型。

%% 【此处粘贴：全市综合典型日曲线 p1_f04】

\textbf{4. 工作日/周末维度}

定义差异比：
\begin{equation}
\eta_{i}=\frac{D_{i}^{\mathrm{we}}-D_{i}^{\mathrm{wd}}}{D_{i}^{\mathrm{wd}}}\times100\%\tag{1-4}
\end{equation}
按$\pm20\%$阈值分类：$\eta_{i}\ge20\%$为周末主导型，$-20\%<\eta_{i}<20\%$为均衡型，$\eta_{i}\le-20\%$为工作日主导型。

\begin{table}[H]
\centering
\caption{工作日—周末差异特征}
\begin{tabular}{ccc}
\toprule
区域 & $\eta_{i}$ & 日类型特征 \\
\midrule
$9$ & $99.63\%$ & 周末主导型 \\
$5$ & $41.89\%$ & 周末主导型 \\
$7$ & $11.74\%$ & 均衡型 \\
$10$ & $-4.87\%$ & 均衡型 \\
$3$ & $-17.08\%$ & 均衡型 \\
$4$ & $-21.00\%$ & 工作日主导型 \\
$1$ & $-24.09\%$ & 工作日主导型 \\
$2$ & $-24.87\%$ & 工作日主导型 \\
$6$ & $-27.35\%$ & 工作日主导型 \\
$8$ & $-39.36\%$ & 工作日主导型 \\
\bottomrule
\end{tabular}
\end{table}

区域$9$周末需求接近工作日$2$倍，后续规划不能仅按工作日配置。

\subsubsection{亮出总模型：三因素岭回归}

候选解释变量含人口密度$\mathrm{PD}_{i}$、车流量$\mathrm{TF}_{i}$、商业POI$\mathrm{POI}_{i}$及既有桩数。剔除存在严格线性关系且具内生性的桩数变量，保留：
\begin{equation}
\mathbf{x}_{i}=(\mathrm{PD}_{i},\mathrm{TF}_{i},\mathrm{POI}_{i})\tag{1-5}
\end{equation}
特征标准化$z_{ij}=(x_{ij}-\mu_{j})/\sigma_{j}$，建立岭回归：
\begin{equation}
\hat{D}_{i}=\beta_{0}+\beta_{1}z_{i,\mathrm{PD}}+\beta_{2}z_{i,\mathrm{TF}}+\beta_{3}z_{i,\mathrm{POI}}\tag{1-6}
\end{equation}
参数通过下式求解：
\begin{equation}
\min_{\beta_{0},\boldsymbol{\beta}}\left[\sum_{i=1}^{10}(D_{i,0}-\beta_{0}-\mathbf{z}_{i}^{\mathrm{T}}\boldsymbol{\beta})^{2}+\lambda\sum_{j=1}^{3}\beta_{j}^{2}\right]\tag{1-7}
\end{equation}
在$\lambda\in\{0.01,0.1,1,10,100\}$上经留一交叉验证以RMSE最小择优，得$\lambda^{*}=10$。

\subsubsection{详述求解}

采用严格LOOCV：每次留$1$区域测试、$9$区域训练，训练折内完成标准化，避免信息泄露。全样本拟合得：
\begin{equation}
\boxed{\hat{D}_{i}=13800.70+513.57z_{i,\mathrm{PD}}+976.24z_{i,\mathrm{TF}}+679.23z_{i,\mathrm{POI}}}\tag{1-8}
\end{equation}
其中$\mu_{\mathrm{PD}}=1996.10,\sigma_{\mathrm{PD}}=1178.14$；$\mu_{\mathrm{TF}}=21265.70,\sigma_{\mathrm{TF}}=7302.37$；$\mu_{\mathrm{POI}}=677.80,\sigma_{\mathrm{POI}}=367.07$。

模型验证如表4：

\begin{table}[H]
\centering
\caption{LOOCV模型对比}
\begin{tabular}{lccccc}
\toprule
模型 & 特征 & RMSE & MAE & MAPE & $R^{2}_{\mathrm{CV}}$ \\
\midrule
均值基准 & 无 & $4336.10$ & $3833.67$ & $34.17\%$ & $-0.2346$ \\
四特征岭回归 & 人口、车流、POI、总桩数 & $3401.18$ & $2967.11$ & $25.97\%$ & $0.2404$ \\
\textbf{三因素岭回归} & \textbf{人口、车流、POI} & $\mathbf{3426.49}$ & $3057.64$ & $26.76\%$ & $0.2291$ \\
对数三因素 & $\ln(x+1)$ & $3738.13$ & $3286.49$ & $28.52\%$ & $0.0825$ \\
低复杂度LightGBM & 四特征 & $4336.10$ & $3833.67$ & $34.17\%$ & $-0.2346$ \\
\bottomrule
\end{tabular}
\end{table}

四特征模型RMSE仅比三因素低$25.31\,\mathrm{kWh/day}$（$0.74\%$），考虑桩数内生性，最终选用三因素模型；LightGBM在每折$9$样本下未形成有效分裂，退化为均值预测。

\subsubsection{结果展示与分析}

\textbf{1. 未来情景预测框架}

一般形式为$\widehat{D}_{i,\tau}=D_{i,0}\cdot f(\mathbf{x}_{i,\tau})/f(\mathbf{x}_{i,0})\cdot(1+g_{\mathrm{EV}})^{\tau}$。短期$\mathbf{x}_{i,1}=\mathbf{x}_{i,0}$，故：
\begin{equation}
\boxed{\widehat{D}_{i,2026}^{(g)}=D_{i,2025}(1+g)}\tag{1-9}
\end{equation}
设置$g\in\{10\%,15\%,20\%\}$三情景：

\begin{table}[H]
\centering
\caption{2026年三情景预测结果（$\mathrm{kWh/day}$）}
\begin{tabular}{ccccc}
\toprule
区域 & 2025基准 & 保守$10\%$ & 基准$15\%$ & 加速$20\%$ \\
\midrule
$1$ & $17884.29$ & $19672.71$ & $20566.93$ & $21461.14$ \\
$2$ & $16504.43$ & $18154.87$ & $18980.09$ & $19805.31$ \\
$3$ & $12828.00$ & $14110.80$ & $14752.20$ & $15393.60$ \\
$4$ & $18694.57$ & $20564.03$ & $21498.76$ & $22433.49$ \\
$5$ & $15076.43$ & $16584.07$ & $17337.89$ & $18091.71$ \\
$6$ & $16940.00$ & $18634.00$ & $19481.00$ & $20328.00$ \\
$7$ & $7274.00$ & $8001.40$ & $8365.10$ & $8728.80$ \\
$8$ & $8452.86$ & $9298.14$ & $9720.79$ & $10143.43$ \\
$9$ & $14956.00$ & $16451.60$ & $17199.40$ & $17947.20$ \\
$10$ & $9396.43$ & $10336.07$ & $10805.89$ & $11275.71$ \\
\midrule
\textbf{全市} & $\mathbf{138007.00}$ & $\mathbf{151807.70}$ & $\mathbf{158708.05}$ & $\mathbf{165608.40}$ \\
\bottomrule
\end{tabular}
\end{table}

\textbf{2. 深入分析}
\begin{itemize}
\item \textbf{现象：}区域需求分层显著，全市峰值集中于晚高峰，区域$9$周末主导特征突出。
\item \textbf{原因：}这主要因为区域功能分异导致人口、交通、商业POI分布不均，且用户出行习惯在工作日通勤与周末休闲场景下差异显著，叠加小样本共线性使单一因素因果难以剥离。
\item \textbf{结论：}三因素岭回归在$10\%$—$20\%$增长率区间提供了稳健的短期预测；$15\%$基准情景全市$158708.05\,\mathrm{kWh/day}$将作为问题二默认输入。
\end{itemize}

%------------------------------------------------
\subsection{问题二：新增快慢充桩多目标整数规划}

\subsubsection{开篇明义}
本节以问题一$15\%$基准车次需求为输入，构建“区域硬约束+全市多目标”规划框架，实现成本、覆盖、接入压力三目标协同优化，经NSGA-II求解与多级筛选确定最优配置。

\subsubsection{分步推导}

\textbf{1. 需求与覆盖参数}

2025加权日均车次$Q_{i,2025}=(5Q_{i}^{\mathrm{wd}}+2Q_{i}^{\mathrm{we}})/7$，规划期需求：
\begin{equation}
Q_{i}^{\mathrm{plan}}=1.15\,Q_{i,2025}\tag{2-1}
\end{equation}
由附件$5$得$c_{f}=6$万元/台，$c_{s}=0.8$万元/台，$p_{f}=120\,\mathrm{kW/台}$，$p_{s}=7\,\mathrm{kW/台}$，$s_{f}=80$，$s_{s}=20$车次/(台$\cdot$日)；附件$1$电网容量$g_{i}$万千瓦换算$G_{i}=10^{4}g_{i}$。

单桩有效覆盖面积反推：
\begin{equation}
a_{i}=A_{i}^{(0)}/N_{i}^{(0)}\tag{2-2}
\end{equation}
新增后区域覆盖率：
\begin{equation}
C_{i}=\min\left\{1,\frac{A_{i}^{(0)}+a_{i}(x_{i}+y_{i})}{A_{i}}\right\}\tag{2-3}
\end{equation}
全市面积加权覆盖率$C_{\mathrm{city}}=\sum A_{i}C_{i}/\sum A_{i}$，既有快充占比$\theta_{i}=F_{i}^{(0)}/(F_{i}^{(0)}+S_{i}^{(0)})$。

\textbf{2. 目标函数}

目标$1$（成本最小）：
\begin{equation}
\min f_{1}=\sum_{i=1}^{10}(6x_{i}+0.8y_{i})\tag{2-4}
\end{equation}
目标$2$（覆盖率最大）：$\max f_{2}=C_{\mathrm{city}}$，NSGA-II中以$-C_{\mathrm{city}}$最小化。

目标$3$（接入压力）：新增接入功率$P_{i}^{\mathrm{add}}=120x_{i}+7y_{i}$，压力率$\nu_{i}=P_{i}^{\mathrm{add}}/G_{i}$。构建两种候选：

$M\text{-}A$压力平方和：
\begin{equation}
\min f_{3}^{(A)}=\sum_{i=1}^{10}\nu_{i}^{2}\tag{2-5}
\end{equation}
$M\text{-}B$压力方差：$\bar{\nu}=\frac{1}{10}\sum\nu_{i}$，
\begin{equation}
\min f_{3}^{(B)}=\frac{1}{10}\sum_{i=1}^{10}(\nu_{i}-\bar{\nu})^{2}\tag{2-6}
\end{equation}
二者关系$\sum\nu_{i}^{2}=10\mathrm{Var}+10\bar{\nu}^{2}$，平方和同时控制总体压力与局部集中。

\subsubsection{亮出总模型}

\begin{equation}
\begin{aligned}
\min\ & [f_{1},\ -C_{\mathrm{city}},\ f_{3}] \\
\mathrm{s.t.}\ & A_{i}^{(0)}+a_{i}(x_{i}+y_{i})\ge0.9A_{i},\quad\forall i \\
& 80(F_{i}^{(0)}+x_{i})+20(S_{i}^{(0)}+y_{i})\ge Q_{i}^{\mathrm{plan}},\quad\forall i \\
& S_{i}^{(0)}x_{i}\ge F_{i}^{(0)}y_{i},\quad\forall i \\
& x_{i},y_{i}\in\mathbb{Z}_{\ge0},\quad 0\le x_{i}+y_{i}\le n_{i}^{100}=\lceil(A_{i}-A_{i}^{(0)})/a_{i}\rceil
\end{aligned}\tag{2-7}
\end{equation}
另设逐时安全预检查$L_{i,t}^{\mathrm{plan}}=1.15E_{i,t}^{*}\le G_{i,t}$，全部区域通过，区域$9$负荷率$0.5882$为最紧。

\subsubsection{详述求解}

\textbf{1. NSGA-II生成Pareto集} 个体编码$\mathbf{z}=(x_{1},y_{1},\ldots,x_{10},y_{10})$共$20$维整数，种群$240$、迭代$500$代（实$417$代收敛），$5$种子重复检验稳定性。

%% 【此处粘贴：NSGA-II收敛过程 p2_f03】

\textbf{2. 最小快充结构修复} 对给定$n_{i}=x_{i}+y_{i}$，最小快充数$x_{i}^{\min}(n_{i})=\lceil\theta_{i}n_{i}\rceil$，$y_{i}^{\min}=n_{i}-x_{i}^{\min}$；若服务能力不足则逐步以快充替换慢充。修复消除区域$3$冗余：$(11,5)\to(10,6)$，成本由$70.0$降至$64.8$万元，接入功率由$1355$降至$1242\,\mathrm{kW}$。

\textbf{3. 邻域支配扫描} 对修复后解搜索$(x_{i}+\Delta x,y_{i}+\Delta y)$，$\Delta x,\Delta y\in\{-1,0,1\}$，检验是否存在严格支配可行邻域点，最终推荐方案通过检验。

\textbf{4. 熵权-TOPSIS择优} 以$[f_{1},C_{\mathrm{city}},f_{3}]$构造评价矩阵，熵权确定客观权重，TOPSIS计算贴近度择优；另以等权与偏成本$(0.50,0.25,0.25)$检验敏感性，M-A三权重成本差异$\le2.2\%$。

%% 【此处粘贴：修复后Pareto前沿二维投影 p2_f04；TOPSIS权重敏感性 p2_f09】

\subsubsection{结果展示与分析}

\textbf{1. M-A与M-B对比择优}

\begin{table}[H]
\centering
\caption{M-A与M-B熵权-TOPSIS推荐方案对比}
\begin{tabular}{lcc}
\toprule
指标 & M-A：压力平方和 & M-B：压力方差 \\
\midrule
总成本/万元 & $\mathbf{4136.80}$ & $4579.20$ \\
面积加权覆盖率 & $90.61\%$ & $\mathbf{92.88\%}$ \\
平均接入压力$\bar{\nu}$ & $\mathbf{0.0367}$ & $0.0396$ \\
最大接入压力$\nu_{\max}$ & $0.1132$ & $0.1132$ \\
压力标准差$\sigma_{\nu}$ & $0.0344$ & $\mathbf{0.0322}$ \\
压力平方和 & $\mathbf{0.0254}$ & $0.0261$ \\
压力方差 & $0.0012$ & $\mathbf{0.0010}$ \\
\bottomrule
\end{tabular}
\end{table}

M-B虽降低标准差，但成本高$10.7\%$且平均压力更高，属“总体压力更高但分布更均匀”的伪均衡，故选用M-A。

%% 【此处粘贴：两类压力目标方案对比 p2_f05】

\textbf{2. 最终配置方案}

\begin{table}[H]
\centering
\caption{最终新增快慢充桩配置（M-A修复后推荐）}
\begin{tabular}{cccccccccc}
\toprule
区域 & 新增快充 & 新增慢充 & 合计 & 成本/万元 & 覆盖率 & 接入功率/kW & $\nu_{i}$ & 服务裕度 \\
\midrule
$1$ & $15$ & $10$ & $25$ & $98.00$ & $90.15\%$ & $1870$ & $0.0058$ & $12338.63$ \\
$2$ & $25$ & $16$ & $41$ & $162.80$ & $94.02\%$ & $3112$ & $0.0104$ & $12286.26$ \\
$3$ & $10$ & $6$ & $16$ & $64.80$ & $90.27\%$ & $1242$ & $0.0045$ & $9168.04$ \\
$4$ & $89$ & $58$ & $147$ & $580.40$ & $90.38\%$ & $11086$ & $0.0315$ & $16900.76$ \\
$5$ & $96$ & $63$ & $159$ & $626.40$ & $91.61\%$ & $11961$ & $0.0532$ & $14775.04$ \\
$6$ & $29$ & $18$ & $47$ & $188.40$ & $90.46\%$ & $3606$ & $0.0117$ & $10383.43$ \\
$7$ & $117$ & $78$ & $195$ & $764.40$ & $90.14\%$ & $14586$ & $0.0784$ & $14427.21$ \\
$8$ & $94$ & $62$ & $156$ & $613.60$ & $90.68\%$ & $11714$ & $0.0459$ & $13188.69$ \\
$9$ & $138$ & $92$ & $230$ & $901.60$ & $90.53\%$ & $17204$ & $0.1132$ & $15195.36$ \\
$10$ & $21$ & $13$ & $34$ & $136.40$ & $90.14\%$ & $2611$ & $0.0127$ & $6032.78$ \\
\midrule
\textbf{全市} & $\mathbf{634}$ & $\mathbf{416}$ & $\mathbf{1050}$ & $\mathbf{4136.80}$ & $\mathbf{90.61\%}$ & $\mathbf{78992}$ & $\mathbf{0.0367}$ & --- \\
\bottomrule
\end{tabular}
\end{table}

快充占比$634/1050=60.38\%$，维持快充主导结构。区域$7,8,9$新增量大主因是面积大、既有覆盖率低（$25.04\%$、$34.99\%$、$19.95\%$），需大量布点达$90\%$底线；区域$9$接入压力最高为后续重点监测对象。

%% 【此处粘贴：区域覆盖率建设前后对比 p2_f01；服务能力约束核验 p2_f02；最终新增分区分布 p2_f06；区域成本与接入压力 p2_f07；约束裕度 p2_f08】

\textbf{3. 深入分析}
\begin{itemize}
\item \textbf{现象：}全市覆盖率达标且成本、压力协同最优，区域间新增量与需求不完全正相关。
\item \textbf{原因：}这主要因为覆盖率硬约束由面积与既有覆盖共同决定，大面积低覆盖区域需更多桩数填补空间盲区，而快充结构约束与压力平方和目标共同抑制了向容量紧张区域过度集中快充。
\item \textbf{结论：}方案满足全部硬约束且通过邻域最优性检验，为后续分时调度提供了坚实的物理网络基础。
\end{itemize}

%------------------------------------------------
\subsection{问题三：基于分时电价的峰谷负荷转移优化}

\subsubsection{开篇明义}
本节直接使用附件$3$、$4$负荷数据，构建峰谷差与方差字典序优化模型，将高峰$20\%$负荷最优分配至低谷，实现削峰填谷与容量安全协同。

\subsubsection{分步推导}

\textbf{1. 时段划分与符号}

题设时段：
\begin{equation}
\begin{aligned}
\mathcal{V}&=\{00\text{--}01,\ldots,06\text{--}07\},\ |\mathcal{V}|=7\\
\mathcal{M}&=\{07\text{--}08,\ldots,10\text{--}11,14\text{--}15,15\text{--}16\},\ |\mathcal{M}|=6\\
\mathcal{H}&=\{11\text{--}12,12\text{--}13,13\text{--}14,16\text{--}17,\ldots,22\text{--}23\},\ |\mathcal{H}|=10
\end{aligned}\tag{3-1}
\end{equation}
$23\text{--}00$保持不变不参与转移。记$L_{i,t}^{d}$为调度前平均负荷，$G_{i,t}$为允许负荷，$z_{i,t}^{d}$为低谷转移量。

\textbf{2. 高峰转移规则}

转移总量：
\begin{equation}
M_{i}^{d}=0.2\sum_{t\in\mathcal{H}}L_{i,t}^{d}\tag{3-2}
\end{equation}
高峰削减：$\widetilde{L}_{i,t}^{d}=0.8L_{i,t}^{d},t\in\mathcal{H}$；平段不变；低谷$\widetilde{L}_{i,t}^{d}=L_{i,t}^{d}+z_{i,t}^{d},t\in\mathcal{V}$，且$\sum_{t\in\mathcal{V}}z_{i,t}^{d}=M_{i}^{d}$，保证$\sum_{t}\widetilde{L}_{i,t}^{d}=\sum_{t}L_{i,t}^{d}$。

\textbf{3. 调度前峰谷差}

\begin{equation}
\Delta_{i}^{d,\mathrm{pre}}=\max_{t}L_{i,t}^{d}-\min_{t}L_{i,t}^{d}\tag{3-3}
\end{equation}

\subsubsection{亮出总模型：两阶段字典序优化}

\textbf{第一阶段：峰谷差最小（LP）}
\begin{equation}
\begin{aligned}
\min\ & U_{i}^{d}-V_{i}^{d}\\
\mathrm{s.t.}\ & V_{i}^{d}\le\widetilde{L}_{i,t}^{d}\le U_{i}^{d},\ \forall t\\
& \sum_{t\in\mathcal{V}}z_{i,t}^{d}=M_{i}^{d},\ 0\le z_{i,t}^{d}\le G_{i,t}-L_{i,t}^{d}\\
& \widetilde{L}_{i,t}^{d}\le G_{i,t},\ \forall t
\end{aligned}\tag{3-4}
\end{equation}
最优值$\Delta_{i}^{d,*}$。

\textbf{第二阶段：方差最小（QP）}
\begin{equation}
\begin{aligned}
\min\ & \frac{1}{24}\sum_{t}(\widetilde{L}_{i,t}^{d}-\bar{L}_{i}^{d})^{2}\\
\mathrm{s.t.}\ & \text{式(3-4)全部约束},\ U_{i}^{d}-V_{i}^{d}\le\Delta_{i}^{d,*}+10^{-6}
\end{aligned}\tag{3-5}
\end{equation}
该模型具水位填充解释：优先抬升原始负荷较低的低谷小时，使低谷曲线平坦化。

\subsubsection{详述求解}

$10\times2=20$个独立子问题。第一阶段\texttt{scipy.optimize.linprog(method='highs')}求解；第二阶段先SLSQP，$5$场景数值回退后以\texttt{CVXPY+CLARABEL}补算，$20$题全部可行。

可行性预检：低谷接纳能力$B_{i}^{d}=\sum_{t\in\mathcal{V}}\max\{0,G_{i,t}-L_{i,t}^{d}\}$均满足$B_{i}^{d}\ge M_{i}^{d}$，最紧为区域$9$周末$B_{9}^{\mathrm{we}}/M_{9}^{\mathrm{we}}\approx4.67$。核验满足电量守恒误差$<10^{-8}$、转移闭合$<10^{-8}$、容量约束$\max(\widetilde{L}-G)\le0$。

%% 【此处粘贴：转移量与低谷接纳能力对比 p3_f02；低谷水位填充示例 p3_f11；LP与QP分配对比 p3_f12】

\subsubsection{结果展示与分析}

\begin{table}[H]
\centering
\caption{工作日调度效果}
\begin{tabular}{ccccccc}
\toprule
区域 & 调度前峰谷差 & 调度后峰谷差 & 改善率 & 方差改善率 & 调度后负荷率 & 安全裕度/kW \\
\midrule
$1$ & $2202.0$ & $1521.8$ & $30.89\%$ & $63.40\%$ & $45.59\%$ & $894.9$ \\
$2$ & $2063.0$ & $1465.4$ & $28.97\%$ & $62.76\%$ & $42.66\%$ & $925.2$ \\
$3$ & $1521.0$ & $1094.2$ & $28.06\%$ & $63.44\%$ & $31.49\%$ & $1065.0$ \\
$4$ & $1969.0$ & $1280.0$ & $\mathbf{34.99\%}$ & $67.47\%$ & $40.75\%$ & $849.7$ \\
$5$ & $1506.0$ & $1021.2$ & $32.19\%$ & $59.84\%$ & $31.08\%$ & $1088.7$ \\
$6$ & $1750.0$ & $1168.6$ & $33.22\%$ & $66.77\%$ & $36.17\%$ & $908.7$ \\
$7$ & $719.0$ & $504.6$ & $29.82\%$ & $65.08\%$ & $14.92\%$ & $1229.6$ \\
$8$ & $980.0$ & $725.0$ & $26.02\%$ & $49.17\%$ & $20.37\%$ & $1207.7$ \\
$9$ & $1210.0$ & $838.4$ & $30.71\%$ & $57.30\%$ & $25.00\%$ & $1138.2$ \\
$10$ & $899.0$ & $603.4$ & $32.88\%$ & $\mathbf{68.10\%}$ & $18.64\%$ & $1150.8$ \\
\midrule
均值 & --- & --- & $\mathbf{30.78\%}$ & $\mathbf{62.33\%}$ & --- & --- \\
\bottomrule
\end{tabular}
\end{table}

全市工作日峰谷差$14057\to9669\,\mathrm{kW}$，改善$\mathbf{31.2\%}$。

\begin{table}[H]
\centering
\caption{周末调度效果}
\begin{tabular}{ccccccc}
\toprule
区域 & 调度前峰谷差 & 调度后峰谷差 & 改善率 & 方差改善率 & 调度后负荷率 & 安全裕度/kW \\
\midrule
$1$ & $1302.0$ & $1068.0$ & $17.97\%$ & $50.06\%$ & $33.46\%$ & $1055.9$ \\
$2$ & $1277.0$ & $1095.0$ & $14.25\%$ & $47.41\%$ & $32.85\%$ & $1106.1$ \\
$3$ & $1157.0$ & $1042.0$ & $9.94\%$ & $44.36\%$ & $29.74\%$ & $1167.7$ \\
$4$ & $1395.0$ & $1090.0$ & $21.86\%$ & $51.70\%$ & $35.82\%$ & $1024.5$ \\
$5$ & $1808.0$ & $1440.0$ & $20.35\%$ & $47.83\%$ & $41.47\%$ & $996.7$ \\
$6$ & $1267.0$ & $1094.0$ & $13.65\%$ & $47.90\%$ & $32.51\%$ & $1103.2$ \\
$7$ & $894.0$ & $729.0$ & $18.46\%$ & $39.22\%$ & $18.71\%$ & $1245.1$ \\
$8$ & $611.0$ & $611.0$ & $0.00\%$ & $41.42\%$ & $15.68\%$ & $1298.6$ \\
$9$ & $2285.0$ & $1853.0$ & $18.91\%$ & $43.08\%$ & $52.30\%$ & $614.0$ \\
$10$ & $952.0$ & $717.0$ & $24.68\%$ & $42.60\%$ & $19.81\%$ & $1205.3$ \\
\midrule
均值 & --- & --- & $\mathbf{16.01\%}$ & $\mathbf{45.56\%}$ & --- & --- \\
\bottomrule
\end{tabular}
\end{table}

全市周末峰谷差$11896\to10484\,\mathrm{kW}$，改善$\mathbf{11.9\%}$。

%% 【此处粘贴：调度前峰谷差对比 p3_f01；全市工作日曲线 p3_f03；全市周末曲线 p3_f04；区域峰谷差改善率 p3_f05；方差改善率 p3_f06；调度后负荷率与安全裕度 p3_f07】

\textbf{典型区域解读：}
\begin{itemize}
\item 区域$4$工作日：峰谷差$1969.0\to1280.0\,\mathrm{kW}$，改善$34.99\%$，方差降$67.47\%$，为峰值与高峰时段高度匹配时的典型削峰效果。
\item 区域$8$周末：峰谷差$0\%$改善（峰值$611\,\mathrm{kW}$位于$10\text{--}11$平段、谷值位于$23\text{--}00$不调度时段），但方差仍降$41.42\%$，揭示固定电价时段与实际峰值错位时的政策边界。
\item 区域$9$周末：$12\text{--}13$时$2285\to1828\,\mathrm{kW}$，调度后全天最大$2038\,\mathrm{kW}$位于$14\text{--}15$平段，最小裕度$614\,\mathrm{kW}$仍安全。
\end{itemize}

%% 【此处粘贴：区域4工作日调度曲线 p3_f08；区域8周末曲线 p3_f09；区域9周末曲线 p3_f10】

\textbf{电网安全：}调度后最大负荷率$\le52.30\%$（区域$9$周末），最小裕度$614.0\,\mathrm{kW}$，无逐时越限；原$3$个超$2100\,\mathrm{kW}$小时（区域$1$工作日$17\text{--}18$时$2221$，区域$9$周末$12\text{--}13$时$2285$及$13\text{--}14$时$2156$）均降至阈值以下。

\textbf{深入分析：}
\begin{itemize}
\item \textbf{现象：}工作日改善显著优于周末，区域$8$周末峰谷差未变但方差下降。
\item \textbf{原因：}这主要因为工作日负荷峰值多位于题设高峰时段，削减直接降低峰值；而周末部分区域峰值位于平段或不调度时段，高峰削减无法触及全天峰值，但低谷填充仍平滑了曲线。
\item \textbf{结论：}分时电价在峰值匹配时效果显著，否则需精细化调整峰平谷时段划分。
\end{itemize}

%------------------------------------------------
\subsection{问题四：三年增长风险评估与滚动优化}

\subsubsection{开篇明义}
本节以问题二已建网络为初始状态、问题三调度曲线为基线，按年增$15\%$投影2026—2028年，分离服务侧与电网侧风险，提出滚动补桩与配网治理策略。

\subsubsection{分步推导}

\textbf{1. 年度需求预测}

$Q_{i,0}=(5Q_{i}^{\mathrm{wd}}+2Q_{i}^{\mathrm{we}})/7$，$Q_{i,k}=Q_{i,0}g_{k}$，其中$g_{1}=1.15$，$g_{2}=1.3225$，$g_{3}=1.520875$。以区域$10$为例：$Q_{10,0}=736.71$，则$Q_{10,2026}=847.22$，$Q_{10,2027}=974.30$，$Q_{10,2028}=1120.45$车次/日。

\textbf{2. 服务侧静态风险}

问题二实施后$F_{i,1}=F_{i}^{(0)}+x_{i}^{(P2)}$，$S_{i,1}=S_{i}^{(0)}+y_{i}^{(P2)}$，服务能力$SC_{i}=80F_{i,1}+20S_{i,1}$。不补建时$SC_{i,k}^{S0}=SC_{i}$，保障率$\psi_{i,k}^{S0}=SC_{i}/Q_{i,k}$，缺口$\mathrm{Gap}_{i,k}^{\mathrm{srv}}=[Q_{i,k}-SC_{i}]_{+}$，按$\psi\ge1$安全、$0.9\le\psi<1$预警、$\psi<0.9$不足分级。

\textbf{3. 电网侧峰时风险投影}

以问题三更危险日期类型的最大负荷率为基线：$\rho_{i,0}^{\mathrm{tou}}=\max_{d}\max_{t}\widetilde{L}_{i,t}^{d}/G_{i,t}$，未来$\rho_{i,k}^{\mathrm{tou}}=\rho_{i,0}^{\mathrm{tou}}g_{k}$，按$\rho<0.80$安全、$0.80\le\rho<1.00$预警、$\rho\ge1.00$过载分级。

\subsubsection{亮出总模型：滚动补桩MILP}

若需追加，设$u_{i,k},v_{i,k}\in\mathbb{Z}_{\ge0}$为2027、2028年追加快、慢充，状态$F_{i,k}=F_{i,k-1}+u_{i,k}$，$S_{i,k}=S_{i,k-1}+v_{i,k}$，约束$80F_{i,k}+20S_{i,k}\ge Q_{i,k}$及快充结构$(1-\theta_{i})u_{i,k}\ge\theta_{i}v_{i,k}$，字典序目标：
\begin{equation}
\min f_{1}=\sum_{i,k}(6u_{i,k}+0.8v_{i,k}),\quad \min f_{2}=\sum_{i}(6u_{i,2}+0.8v_{i,2})\tag{4-1}
\end{equation}
即先总成本最小，再2027年投资最小。

另对过载区域计算最低技术容量缺口：以区域$9$控制时段$10\text{--}11$时为例，$\widetilde{L}_{9,10}=1915$，$G_{9,10}=2027\,\mathrm{kW}$，
\begin{equation}
\Delta G_{9,k}^{\min}=[1915g_{k}-2027]_{+}\tag{4-2}
\end{equation}

\subsubsection{详述求解}
采用混合整数线性规划求解滚动模型。

\subsubsection{结果展示与分析}

\textbf{1. 服务能力风险：无缺口}

\begin{table}[H]
\centering
\caption{静态服务保障率（$S0$不补建）}
\begin{tabular}{ccccccc}
\toprule
区域 & $SC_{i}$/车次 & 2026$\psi$ & 2027$\psi$ & 2028$\psi$ & 最低$\psi$ & 等级 \\
\midrule
$1$ & $13440$ & $1220.4\%$ & $1061.1\%$ & $922.7\%$ & $922.7\%$ & 安全 \\
$2$ & $13420$ & $1183.7\%$ & $1029.3\%$ & $895.0\%$ & $895.0\%$ & 安全 \\
$3$ & $10160$ & $1024.3\%$ & $890.8\%$ & $774.5\%$ & $774.5\%$ & 安全 \\
$4$ & $18460$ & $1183.9\%$ & $1029.5\%$ & $895.2\%$ & $895.2\%$ & 安全 \\
$5$ & $16020$ & $1286.8\%$ & $1119.0\%$ & $973.0\%$ & $973.0\%$ & 安全 \\
$6$ & $11540$ & $997.8\%$ & $867.6\%$ & $754.5\%$ & $754.5\%$ & 安全 \\
$7$ & $15120$ & $2182.5\%$ & $1897.9\%$ & $1650.3\%$ & $1650.3\%$ & 安全 \\
$8$ & $14260$ & $1331.0\%$ & $1157.5\%$ & $1006.5\%$ & $1006.5\%$ & 安全 \\
$9$ & $16520$ & $1247.2\%$ & $1084.5\%$ & $943.0\%$ & $943.0\%$ & 安全 \\
$10$ & $6880$ & $812.1\%$ & $706.2\%$ & $\mathbf{614.0\%}$ & $\mathbf{614.0\%}$ & 安全 \\
\bottomrule
\end{tabular}
\end{table}

至2028年$\min\psi_{i,2028}=614.0\%>100\%$，$\mathrm{Gap}_{i,k}^{\mathrm{srv}}=0$，地理覆盖率保持问题二水平不因车辆增长而下降。

%% 【此处粘贴：S0服务保障率年度变化 p4_f01；全市服务保障率年度对比 p4_f04；滚动补桩前后能力对比 p4_f03】

\textbf{2. 峰时电网风险：区域9过载}

\begin{table}[H]
\centering
\caption{基期最大负荷率与TOU缓解}
\begin{tabular}{ccccc}
\toprule
区域 & 调度前负荷率 & 调度后负荷率 & 缓解量 & 控制日类型 \\
\midrule
$1$ & $43.99\%$ & $35.19\%$ & $8.80\%$ & 工作日 \\
$2$ & $42.42\%$ & $33.93\%$ & $8.48\%$ & 工作日 \\
$3$ & $33.37\%$ & $32.47\%$ & $0.90\%$ & 周末 \\
$4$ & $35.47\%$ & $30.57\%$ & $4.89\%$ & 工作日 \\
$5$ & $52.58\%$ & $52.58\%$ & $0.00\%$ & 工作日 \\
$6$ & $36.56\%$ & $31.14\%$ & $5.43\%$ & 工作日 \\
$7$ & $30.31\%$ & $30.31\%$ & $0.00\%$ & 工作日 \\
$8$ & $26.32\%$ & $26.32\%$ & $0.00\%$ & 工作日 \\
$\mathbf{9}$ & $\mathbf{94.47\%}$ & $\mathbf{94.47\%}$ & $\mathbf{0.00\%}$ & \textbf{周末} \\
$10$ & $28.31\%$ & $28.31\%$ & $0.00\%$ & 工作日 \\
\bottomrule
\end{tabular}
\end{table}

区域$9$调度后峰时风险率投影：

\begin{table}[H]
\centering
\caption{区域9未来峰时风险率与容量缺口}
\begin{tabular}{ccccc}
\toprule
年份 & $g_{k}$ & 调度后峰时风险率 & 等级 & $\Delta G_{9,k}^{\min}$ \\
\midrule
2026 & $1.15$ & $108.65\%$ & 过载 & $0.175\,\mathrm{MW}$ \\
2027 & $1.3225$ & $124.94\%$ & 过载 & $0.506\,\mathrm{MW}$ \\
2028 & $1.520875$ & $143.68\%$ & 过载 & $0.885\,\mathrm{MW}$ \\
\bottomrule
\end{tabular}
\end{table}

区域$9$瓶颈位于$10\text{--}11$时平段（$1915\,\mathrm{kW}$，允许$2027\,\mathrm{kW}$），该时段未纳入当前高峰电价，故分时调度未缓解其风险；区域$5$在2028年达$79.97\%$接近预警线，其余区域安全。

%% 【此处粘贴：调度后峰时风险率投影 p4_f02】

\textbf{3. 滚动优化与调整建议}

MILP最优解$u_{i,k}^{*}=0,v_{i,k}^{*}=0,\ \forall i,k=2,3$，追加成本$0$万元，表明问题二方案已覆盖三年服务需求，继续补桩无法解决区域$9$的电网瓶颈。

据此提出：
\begin{enumerate}[label=\arabic*.]
\item \textbf{区域9电网重点治理：}立即对$10\text{--}11$时平段开展逐时负荷与真实基础负荷复核，分期增容$0.175/0.506/0.885\,\mathrm{MW}$；
\item \textbf{扩展需求响应：}将该平段纳入预约充电、弹性电价或局部需求响应；
\item \textbf{区域5关注监测：}若实际增速超$15\%$或形态恶化，提前启动逐时复核；
\item \textbf{其余区域常规监测：}暂不新增；
\item \textbf{不以补桩替代电网措施：}区域$9$服务保障率2028年仍$943.0\%$，追加桩数反增接入压力。
\end{enumerate}

\textbf{深入分析：}
\begin{itemize}
\item \textbf{现象：}服务侧全安全、电网侧单点过载的分离格局。
\item \textbf{原因：}这主要因为附件单桩服务能力参数较高使服务能力冗余充足，而电网允许容量在区域$9$紧约束且瓶颈时段不在调控范围内，需求同比放大直接触发容量越限。
\item \textbf{结论：}中长期优化重点应由“补桩”转向“平段需求响应与配网增容”。
\end{itemize}

%================================================
\section{模型评价与推广}
\subsection{模型优点}
\begin{enumerate}[label=\arabic*.]
\item 递进式框架完整覆盖需求刻画、设施配置、运行调度、中长期风险四阶段，模型间数据接口清晰，符合城市能源系统规划逻辑。
\item 三因素岭回归经LOOCV与多模型对照验证，在小样本、强共线性条件下实现稳健预测，并以真实观测为基准避免拟合误差叠加。
\item 多目标整数规划设置区域硬约束保障公平性，NSGA-II输出Pareto集避免主观权重，经结构修复与邻域扫描消除算法随机劣解，熵权-TOPSIS提供客观择优。
\item 两阶段负荷转移兼顾峰谷差与方差，显式约束电量守恒与逐时容量，具水位填充可解释性，并揭示固定电价时段的适用边界。
\item 服务—电网分离投影与滚动MILP实现“补桩 vs. 增容”精准分流，避免无效投资。
\end{enumerate}

\subsection{模型局限性与改进}
\begin{enumerate}[label=\arabic*.]
\item 样本仅$10$区域，岭回归系数不能作因果解释；缺少多年序列无法独立识别保有量时间趋势，未来可引入多城市面板数据或贝叶斯分层模型。
\item 覆盖模型以既有网络平均效率近似，未纳入候选站址、道路拓扑及土地可用性；后续可结合GIS与选址集合覆盖模型精细化。
\item $\nu_{i}$为额定接入压力率，未引入同时率与利用率，不能替代潮流计算；后续可耦合配网潮流与利用率情景。
\item 分时调度采用小时级数据，无法刻画$15$分钟瞬时波动；且固定$20\%$转移比例未建模用户价格弹性与停车时长异质性，可引入离散选择模型。
\item $23\text{--}00$时段未分类、跨区域配网拓扑未建模，中长期投影假设形态稳定；后续若获精细数据应滚动校正。
\end{enumerate}

\subsection{灵敏度分析}
\begin{itemize}
\item 需求侧：$10\%$、$15\%$、$20\%$三情景全市需求$151807.70$、$158708.05$、$165608.40\,\mathrm{kWh/day}$，跨度$8.7\%$，区域排序不变，表明规划对增速在$\pm5\%$区间内稳健。
\item 配置侧：M-A方案在熵权、等权、偏成本三权重下成本差异$\le2.2\%$，平均压力差异$\le0.0006$，Pareto筛选稳健。
\item 调度侧：$20\%$转移比例下全部场景接纳能力充足（最小$4.67$倍），但区域$8$周末揭示时段划分敏感性，提示需按实际峰值动态优化峰平谷边界。
\item 风险侧：区域$9$过载阈值对增长率敏感，若增速升至$20\%$，2026年峰时风险率将达$113.37\%$，需更早增容。
\end{itemize}

\subsection{模型推广}
本模型体系可推广至更一般化的城市基础设施智能规划场景：三因素岭回归框架可替换为多源融合的机器学习需求预测；多目标整数规划可扩展至含储能、光储充一体化的选址定容；两阶段负荷转移可推广至考虑用户异质性与动态电价的需求响应；滚动风险评估可接入城市大脑实时数据，实现“预测—配置—调度—增容”闭环的数字孪生城市能源管理平台。

%================================================
\section{参考文献}
\begin{thebibliography}{99}
\bibitem{ref1} H. W. Dommel, W. F. Tinney. Optimal Power Flow Solutions[J]. IEEE Transactions on Power Apparatus and Systems, 1968, 87(10): 1866-1876.
\bibitem{ref2} K. Deb, A. Pratap, S. Agarwal, et al. A Fast and Elitist Multiobjective Genetic Algorithm: NSGA-II[J]. IEEE Transactions on Evolutionary Computation, 2002, 6(2): 182-197.
\bibitem{ref3} A. E. Hoerl, R. W. Kennard. Ridge Regression: Biased Estimation for Nonorthogonal Problems[J]. Technometrics, 1970, 12(1): 55-67.
\bibitem{ref4} 国家市场监督管理总局, 国家标准化管理委员会. GB/T 7714—2015 信息与文献 参考文献著录规则[S]. 北京: 中国标准出版社, 2015.
\bibitem{ref5} 刘文霞, 张建华, 等. 电动汽车充电负荷时空分布预测与配电网协同规划[J]. 电力系统自动化, 2020, 44(12): 45-53.
\bibitem{ref6} 王成山, 王赛一, 等. 考虑分时电价的电动汽车有序充电调度策略[J]. 电网技术, 2019, 43(8): 2756-2764.
\bibitem{ref7} C. L. Hwang, K. Yoon. Multiple Attribute Decision Making: Methods and Applications[M]. Berlin: Springer, 1981.
\bibitem{ref8} S. Shao, M. Pipattanasomporn, S. Rahman. Grid Integration of Electric Vehicles and Demand Response with Customer Choice[J]. IEEE Transactions on Smart Grid, 2012, 3(1): 543-550.
\end{thebibliography}

%================================================
\section{附录}
\subsection{附录A：关键数据表}
\begin{itemize}
\item 表A-1：附件1—5原始数据清洗后汇总表（略，见支撑材料\texttt{problem1\_summary.csv}等）
\item 表A-2：20个区域$\times$日期类型子问题低谷分配$z_{i,t}^{d}$详表（见\texttt{problem3\_summary.csv}）
\end{itemize}

\subsection{附录B：核心代码（节选）}
\begin{lstlisting}[language=Python,caption={\texttt{p1\_ridge.py} 三因素岭回归与LOOCV}]
import numpy as np
from sklearn.linear_model import Ridge
from sklearn.model_selection import LeaveOneOut
# 特征标准化在训练折内完成，避免泄露
for lam in [0.01,0.1,1,10,100]:
    rmse = []
    for train_idx,test_idx in LeaveOneOut().split(X):
        mu,sigma = X[train_idx].mean(0), X[train_idx].std(0,ddof=0)
        Z_train = (X[train_idx]-mu)/sigma
        Z_test = (X[test_idx]-mu)/sigma
        model = Ridge(alpha=lam).fit(Z_train, y[train_idx])
        y_pred = model.predict(Z_test)
        rmse.append((y_pred-y[test_idx])**2)
    print(lam, np.sqrt(np.mean(rmse)))
\end{lstlisting}

\begin{lstlisting}[language=Python,caption={\texttt{p2\_nsga2.py} NSGA-II与结构修复}]
# 个体编码 20维整数 (x1,y1,...,x10,y10)
# 修复：x_min = ceil(theta * n)
def repair(ind):
    for i in range(10):
        n = ind[2*i]+ind[2*i+1]
        x_min = int(np.ceil(theta[i]*n)) if n>0 else 0
        # 保证服务能力与结构约束
        while not feasible(ind): 
            # 慢充替换为快充或增加总数
            pass
    return ind
\end{lstlisting}

\begin{lstlisting}[language=Python,caption={\texttt{p3\_solvers\_lp\_qp.py} 两阶段求解}]
from scipy.optimize import linprog
import cvxpy as cp
# Stage1: linprog HiGHS
res1 = linprog(c, A_ub=A, b_ub=b, A_eq=A_eq, b_eq=b_eq, method='highs')
Delta_star = res1.fun
# Stage2: QP 最小化方差，保留Delta_star约束
x = cp.Variable(7)
objective = cp.Minimize(cp.sum_squares(L_tilde - L_mean))
constraints = [x>=0, cp.sum(x)==M, x<=G_V-L_V, U-V<=Delta_star+1e-6]
prob = cp.Problem(objective, constraints)
prob.solve(solver=cp.CLARABEL)
\end{lstlisting}

完整代码见支撑材料目录\texttt{code/}与\texttt{支撑材料/}，共计$22$个Python模块。

\subsection{附录C：图件索引}
全部图件已输出至\texttt{figures/}目录：
\begin{itemize}
\item 问题一：\texttt{p1\_figures/p1\_f02\_工作日周末全天需求对比.png}、\texttt{p1\_f04\_全市工作日周末综合典型日曲线.png}
\item 问题二：\texttt{p2\_f01}—\texttt{p2\_f09}共$9$幅（覆盖率、服务能力、收敛、Pareto、对比、分布、压力、裕度、敏感性）
\item 问题三：\texttt{p3\_f01}—\texttt{p3\_f12}共$12$幅（峰谷差、接纳能力、全市曲线、改善率、负荷率、区域曲线、水位填充等）
\item 问题四：\texttt{p4\_f01}—\texttt{p4\_f04}共$4$幅（服务保障率、风险投影、滚动对比、年度对比）
\end{itemize}

\end{document}
